%% kapa: exam=1 round=26 year=2015 period=2 form=A subject=appraisal_law subject_ko=부동산관계법규 seq=1 number=1
%% source: Q-Net 감정평가사 기출문제 (cst00309) artlSeq=5138001 file=extracted/1008/5138001/제26회_감정평가사 1차_기출/제26회_감정평가사 1차_기출/2교시_감평_A형.pdf
%% conversion: kordoc 4.24.51 -> tools/gen1.py (auto); review=proofread-sonnet; flags=none
\begin{question}{1}
국토의 계획 및 이용에 관한 법령상의 내용으로 옳지 않은 것은?
\begin{choices}
  \item 기반시설에 관한 사항에 대한 정책방향은 도시ㆍ군기본계획의 내용에 포함되어야 한다.
  \item 입지규제최소구역의 지정 목적을 이루기 위하여 간선도로 등 주요 기반시설의 확보에 관한 사항은 입지규제최소구역계획에 포함되어야 한다.
  \item 지상이나 지하에 기반시설을 설치하려면 원칙적으로 그 시설의 종류ㆍ명칭ㆍ위치ㆍ규모 등을 미리 도시ㆍ군기본계획으로 결정하여야 한다.
  \item 녹지는 기반시설 중 공간시설에 해당한다.
  \item 주거ㆍ상업 또는 공업지역에서의 개발행위로 기반시설의 처리ㆍ공급 또는 수용능력이 부족할 것으로 예상되는 지역 중 기반시설의 설치가 곤란한 지역을 개발밀도관리구역으로 지정할 수 있다.
\end{choices}
\end{question}
\answer{3}
